% Part: lambda-calculus
% Chapter: lambda-definability
% Section: introduction

\documentclass[../../../include/open-logic-section]{subfiles}

\begin{document}

\olfileid{lam}{ldf}{int}
\olsection{مقدمة}

يبدو حساب لامبدا للوهلة الأولى مجرد حساب شديد التجريد لعبارات تمثل
دوال وتطبيقات لهذه الدوال في غيرها. ولا شيء في تركيب حساب لامبدا يوحي
بأنها دوال في أنواع مخصوصة من الأشياء؛ فلا يتضمن التركيب، على وجه
الخصوص، أي ذكر للأعداد الطبيعية. وعمليتاه الأساسيتان---التطبيق وتجريد
لامبدا---عمليتان تنطبقان على أي دالة، لا على الدوال في الأعداد
الطبيعية وحدها.

ومع ذلك، يمكن بشيء من الابتكار تعريف الدوال الحسابية، أي الدوال في
الأعداد الطبيعية، داخل حساب لامبدا. ولعمل ذلك نعرّف، لكل عدد طبيعي
$n \in \Nat$، حد~$\lambd$ مخصوصًا~$\num n$ يسمى \emph{عدد تشيرش}
الموافق لـ$n$. (وقد سميت أعداد تشيرش باسم ألونزو تشيرش.)

\begin{defn}
  إذا كان~$n \in \Nat$، فإن \emph{عدد تشيرش} الموافق~$\num{n}$
  يمثل~$n$:
  \[
    \num{n} \ident \lambd[fx][f^n(x)]
  \]
  ويمثل~$f^n(x)$ هنا نتيجة تطبيق~$f$ في~$x$ عدد~$n$ من المرات.
  فمثلًا، $\num{0}$ هو~$\lambd[fx][x]$، و$\num{3}$ هو
  $\lambd[fx][f(f(f\,x))]$.
\end{defn}
  
يشفر عدد تشيرش~$\num n$ بوصفه حد لامبدا يمثل دالة تقبل وسيطين~$f$
و$x$، وتعيد~$f^n(x)$. ومن الواضح أن أعداد تشيرش في صورة سوية.

ولا يفيد تمثيل الأعداد الطبيعية في حساب لامبدا إلا إذا أمكننا الحساب
بها. والحساب بأعداد تشيرش في حساب لامبدا يعني تطبيق حد~$\lambd$، وليكن
$F$، في عدد تشيرش، واختزال الحد المركب~$F\, \num n$ إلى صورة سوية.
فإذا كان يختزل دائمًا إلى صورة سوية، وكانت هذه الصورة دائمًا عدد
تشيرش~$\num m$، أمكننا اعتبار مخرج الحساب العدد~$m$. وعندئذ يمكننا
اعتبار~$F$ معرّفًا لدالة~$f\colon \Nat \to \Nat$، هي الدالة التي
تحقق~$f(n) = m$ إذا وفقط إذا كان~$F\, \num n \red \num m$. وبفضل
خاصية تشيرش--روسر، تكون الصور السوية وحيدة إن وجدت. لذلك إذا كان
$F\, \num n \red \num m$، فلا يوجد حد آخر في صورة سوية، وبخاصة عدد
تشيرش آخر، يختزل إليه~$F \, \num n$.

وبالعكس، إذا أعطيت دالة~$f\colon \Nat \to \Nat$، أمكننا أن نسأل هل
يوجد حد~$F$ يعرّف~$f$ بهذه الطريقة. وفي هذه الحال نقول إن~$F$
\emph{!!{lambda define}s}~$f$، وإن~$f$ !!{lambda definable}. ويمكن
تعميم ذلك إلى الدوال متعددة المواضع والدوال الجزئية.

\begin{defn}
افترض أن~$f\colon \Nat^k \pto \Nat$ دالة جزئية ذات~$k$ مواضع. نقول
إن حد لامبدا~$F$ \emph{!!{lambda define}s}~$f$ إذا كان، لكل
  $n_0$، \dots،~$n_{k-1}$،
\[
F \, \num{n_0} \, \num{n_1} \dots \num{n_{k-1}} \red \num{f(n_0, n_1, \dots,
  n_{k-1})}
\]
متى كانت~$f(n_0, \dots, n_{k-1})$ معرّفة، ولم يكن للحد
$F \, \num{n_0} \, \num{n_1} \dots \num{n_{k-1}}$ صورة سوية في
غير ذلك.
\end{defn}

ومن أبسط الأمثلة الدوال الثابتة. فالحد~$C_k \ident
\lambd[x][\num{k}]$ !!{lambda define}s الدالة~$c_k\colon \Nat \to
\Nat$ التي تحقق~$c_k(n) = k$. إذ إن
$C_k \, \num n \ident (\lambd[x][\num{k}])\num n \redone \num{k}$
لكل~$n$. والدالة المحايدة !!{lambda defined} بالحد~$\lambd[x][x]$.
أما الدوال الأعقد فأصعب تعريفًا بطبيعة الحال، وغالبًا ما تتطلب قدرًا
كبيرًا من الابتكار. ولذلك قد يكون مدهشًا أن كل دالة قابلة للحساب
!!{lambda definable}. والعكس صحيح أيضًا: إذا كانت دالة ما
!!{lambda definable} فهي قابلة للحساب.

\end{document}
